% =====================================================================
% 6. Requirements and Scope
% =====================================================================
\section{Requirements and Scope}
\label{sec:requirements}

The requirements were derived from the partner's project brief.
Each goal was decomposed into testable requirements linked to a research question, which will be validated with the STT before implementation.
The complete specification of 11 functional and 8 non-functional requirements is provided in Appendix~\ref{sec:trace}.

Three personas were defined.
The \textbf{test engineer} requires visibility of every cell, raw and filtered signals side by side, and early warning of anomalies.
The \textbf{project lead} requires a clear pass or fail outcome supported by the key figures.
The \textbf{customer} encounters the system during a laboratory visit, where it is presented by an engineer.

The essential functional requirements are interchangeable data sources (FR1), KPIs such as pack voltage and cell-voltage spread (FR2), reuse of the STT's calibration code (FR3), near-real-time streaming (FR4), filtering with raw data always retained (FR5), the alert engine (FR6), three role-specific views (FR7) and demonstration replay (FR8).
The principal non-functional requirement is that the dashboard never writes to the rig and imposes negligible load on it (NFR1).
Further targets include an end-to-end latency below 500\,ms (NFR2) and compliance with WCAG~2.2 AA (NFR5).

The project does not modify test procedures or hardware, control the rig, employ machine learning or involve human participants.
Deployment and training constitute company work and fall outside the scope of the dissertation.
